\documentclass[11pt,a4paper]{article}
\usepackage[utf8]{inputenc}
\usepackage[T5]{fontenc}
\usepackage[vietnamese]{babel}
\usepackage[margin=1.4cm, top=1.8cm, bottom=1.8cm]{geometry}
\usepackage{xcolor}
\usepackage{listings}
\usepackage{booktabs}
\usepackage{array}
\usepackage{tabularx}
\usepackage{amsmath}
\usepackage{fancyhdr}
\usepackage{titlesec}
\usepackage{microtype}
\usepackage{enumitem}
\usepackage{hyperref}

\setlength{\emergencystretch}{3em}
\setlist[itemize]{leftmargin=1.2em, itemsep=0.25em, topsep=0.2em}
\urlstyle{tt}

% Màu sắc giao diện báo cáo
\definecolor{PrimaryBlue}{HTML}{1E3A8A}
\definecolor{AccentTeal}{HTML}{0F766E}
\definecolor{CodeBg}{HTML}{F8FAFC}
\definecolor{CodeBorder}{HTML}{CBD5E1}
\definecolor{MutedText}{HTML}{475569}

\hypersetup{
    colorlinks=true,
    linkcolor=PrimaryBlue,
    urlcolor=AccentTeal
}

% Header & Footer
\pagestyle{fancy}
\fancyhf{}
\lhead{\small\color{MutedText} Báo cáo Kỹ thuật: Môi trường Mô phỏng \& Linh kiện Drone X500 HITL}
\rhead{\small\color{MutedText} \thepage}
\renewcommand{\headrulewidth}{0.4pt}

% Định dạng tiêu đề mục
\titleformat{\section}{\Large\bfseries\color{PrimaryBlue}}{\thesection.}{0.5em}{}[\titlerule]
\titleformat{\subsection}{\large\bfseries\color{AccentTeal}}{\thesubsection}{0.5em}{}

% Cấu hình hiển thị code XML/SDF và C++
\lstdefinelanguage{SDF}{
    morekeywords={sdf,world,physics,ode,solver,constraints,plugin,wind,linear_velocity,horizontal,magnitude,sin,direction,gravity,magnetic_field,atmosphere,spherical_coordinates,include,model,link,inertial,mass,inertia,sensor,imu,angular_velocity,linear_acceleration,noise,navsat,position_sensing,velocity_sensing,air_pressure,pressure,magnetometer,joint,axis,limit,dynamics},
    sensitive=true,
    morecomment=[s]{<!--}{-->},
    morestring=[b]",
    morestring=[b]'
}

\lstset{
    basicstyle=\ttfamily\scriptsize,
    backgroundcolor=\color{CodeBg},
    frame=single,
    rulecolor=\color{CodeBorder},
    breaklines=true,
    breakatwhitespace=false,
    showstringspaces=false,
    keywordstyle=\color{PrimaryBlue}\bfseries,
    commentstyle=\color{MutedText}\itshape,
    stringstyle=\color{AccentTeal},
    numbers=left,
    numberstyle=\tiny\color{MutedText},
    numbersep=5pt,
    xleftmargin=12pt,
    framexleftmargin=12pt,
    aboveskip=0pt,
    belowskip=0pt
}

% Macro bố cục chuẩn: Bên trái giải thích (45%) - Bên phải trích xuất code (52%)
\newcommand{\ExplainTitle}[1]{\par\medskip\noindent{\bfseries\color{PrimaryBlue}#1}\par\smallskip}

\begin{document}

\begin{center}
    {\LARGE\bfseries\color{PrimaryBlue} BÁO CÁO MÔI TRƯỜNG MÔ PHỎNG VÀ CẤU HÌNH LINH KIỆN DRONE X500 HITL}\\[0.4em]
    {\normalsize\color{MutedText} Trình bày song song: Giải thích kỹ thuật (bên trái) và Trích xuất mã nguồn cấu hình (bên phải)}
\end{center}
\vspace{0.5em}
\hrule
\vspace{1em}

% ==============================================================================
\section{Môi trường mô phỏng Gazebo Harmonic (\texttt{drone\_world.sdf})}
% ==============================================================================

Môi trường mô phỏng được xây dựng trên nền tảng \textbf{Gazebo Harmonic} thông qua tệp cấu hình \path{simulation/worlds/drone_world.sdf}. Đây là nơi định nghĩa đồng hồ vật lý, các điều kiện tự nhiên (trọng lực, từ trường, gió), tọa độ địa lý gốc và các \emph{System Plugin} phục vụ xuất dữ liệu cảm biến.

% --- 1.1 Physics ---
\subsection{Cấu hình bộ giải vật lý ODE và nhịp mô phỏng}
\noindent
\begin{minipage}[t]{0.45\textwidth}
\vspace{0pt}
\textbf{Giải thích cấu hình:}
\begin{itemize}
    \item \textbf{Bộ giải vật lý:} Sử dụng engine \texttt{ODE} (Open Dynamics Engine) với bộ giải nhanh \texttt{quick} và \texttt{25} vòng lặp/step (\texttt{sor = 1.3}) để cân bằng giữa độ chính xác và tốc độ thời gian thực.
    \item \textbf{Tần số vật lý 250 Hz:} \path{max_step_size = 0.004 s} kết hợp \path{real_time_update_rate = 250} tạo ra nhịp đồng hồ gốc $4\text{ ms/step}$ (\path{real_time_factor = 1.0}).
    \item \textbf{Ý nghĩa với cảm biến:} Mọi \emph{System Plugin} cảm biến đều dựa trên nhịp $250\text{ Hz}$ này để kiểm tra chu kỳ \path{<update_rate>} và phát gói tin lên GZ-Transport.
\end{itemize}
\end{minipage}
\hfill
\begin{minipage}[t]{0.52\textwidth}
\vspace{0pt}
\begin{lstlisting}[language=SDF, firstnumber=5]
<physics type="ode">
  <max_step_size>0.004</max_step_size>
  <real_time_factor>1.0</real_time_factor>
  <real_time_update_rate>250</real_time_update_rate>
  <ode>
    <solver>
      <type>quick</type>
      <iters>25</iters>
      <sor>1.3</sor>
    </solver>
    <constraints>
      <cfm>0.00001</cfm>
      <erp>0.2</erp>
      <contact_max_correcting_vel>100.0</contact_max_correcting_vel>
      <contact_surface_layer>0.001</contact_surface_layer>
    </constraints>
  </ode>
</physics>
\end{lstlisting}
\end{minipage}

\vspace{1em}

% --- 1.2 Environment & GPS Origin ---
\subsection{Điều kiện tự nhiên, tọa độ địa lý gốc và nạp System Plugins}
\noindent
\begin{minipage}[t]{0.45\textwidth}
\vspace{0pt}
\textbf{Giải thích cấu hình:}
\begin{itemize}
    \item \textbf{Trọng lực \& Khí quyển:} Gia tốc trọng trường $g = 9.8\text{ m/s}^2$ hướng xuống (\texttt{-Z}); mô hình khí quyển \texttt{adiabatic} giúp Barometer tính áp suất theo độ cao chính xác.
    \item \textbf{Từ trường Trái Đất:} Vector từ trường ENU $(6\times 10^{-6},\, 2.3\times 10^{-5},\, -4.2\times 10^{-5})\text{ T}$ làm chuẩn cho cảm biến la bàn.
    \item \textbf{Tọa độ GPS gốc:} Chuẩn \texttt{EARTH\_WGS84} hướng \texttt{ENU} tại $(47.397971^\circ,\, 8.546164^\circ,\, 0\text{ m})$ — khớp với gốc \texttt{StateEstimator} trên PC.
    \item \textbf{4 Plugin cảm biến:} Chỉ khi nạp 4 plugin \texttt{Imu}, \texttt{NavSat}, \texttt{AirPressure}, \texttt{Magnetometer} thì Gazebo mới kích hoạt và publish dữ liệu từ các linh kiện gắn trên drone.
\end{itemize}
\end{minipage}
\hfill
\begin{minipage}[t]{0.52\textwidth}
\vspace{0pt}
\begin{lstlisting}[language=SDF, firstnumber=23]
<gravity>0 0 -9.8</gravity>
<magnetic_field>6e-06 2.3e-05 -4.2e-05</magnetic_field>
<atmosphere type="adiabatic"/>

<plugin filename="gz-sim-physics-system" name="gz::sim::systems::Physics"/>
<plugin filename="gz-sim-scene-broadcaster-system" name="gz::sim::systems::SceneBroadcaster"/>
<plugin filename="gz-sim-user-commands-system" name="gz::sim::systems::UserCommands"/>
<plugin filename="gz-sim-imu-system" name="gz::sim::systems::Imu"/>
<plugin filename="gz-sim-navsat-system" name="gz::sim::systems::NavSat"/>
<plugin filename="gz-sim-air-pressure-system" name="gz::sim::systems::AirPressure"/>
<plugin filename="gz-sim-magnetometer-system" name="gz::sim::systems::Magnetometer"/>
...
<spherical_coordinates>
  <surface_model>EARTH_WGS84</surface_model>
  <world_frame_orientation>ENU</world_frame_orientation>
  <latitude_deg>47.397971057728974</latitude_deg>
  <longitude_deg>8.546163739800146</longitude_deg>
  <elevation>0</elevation>
</spherical_coordinates>
\end{lstlisting}
\end{minipage}

\vspace{1em}

% --- 1.3 Wind ---
\subsection{Mô hình nhiễu gió động thực tế (\texttt{WindEffects}) và vị trí khởi tạo Drone}
\noindent
\begin{minipage}[t]{0.45\textwidth}
\vspace{0pt}
\textbf{Giải thích cấu hình:}
\begin{itemize}
    \item \textbf{Gió nền cố định:} Thổi liên tục với vận tốc $1.5\text{ m/s}$ theo trục $X$ và $0.5\text{ m/s}$ theo trục $Y$.
    \item \textbf{Gió giật hình sin (\texttt{WindEffects}):} Biên độ cường độ biến thiên $\pm 30\%$ theo chu kỳ $6.0\text{ s}$ (\texttt{time\_for\_rise = 4.0 s}); hướng gió dao động $\pm 0.3\text{ rad}$ theo chu kỳ $10.0\text{ s}$. Đây là nguồn nhiễu ngoại lực trực tiếp để kiểm chứng khả năng kháng nhiễu của ITSMC và bộ quan sát ESO.
    \item \textbf{Đặt mẫu Drone:} Nạp mô hình \texttt{model://x500} tại tâm tọa độ $(0, 0, 0.12\text{ m})$ trên mặt đất.
\end{itemize}
\end{minipage}
\hfill
\begin{minipage}[t]{0.52\textwidth}
\vspace{0pt}
\begin{lstlisting}[language=SDF, firstnumber=38]
<wind>
  <linear_velocity>1.5 0.5 0</linear_velocity>
</wind>

<plugin filename="gz-sim-wind-effects-system" name="gz::sim::systems::WindEffects">
  <force_approximation_scaling_factor>0.005</force_approximation_scaling_factor>
  <horizontal>
    <magnitude>
      <time_for_rise>4.0</time_for_rise>
      <sin>
        <amplitude_percent>0.30</amplitude_percent>
        <period>6.0</period>
      </sin>
    </magnitude>
    <direction>
      <time_for_rise>8.0</time_for_rise>
      <sin>
        <amplitude>0.3</amplitude>
        <period>10.0</period>
      </sin>
    </direction>
  </horizontal>
</plugin>
...
<include>
  <uri>model://x500</uri>
  <pose>0 0 0.12 0 0 0</pose>
</include>
\end{lstlisting}
\end{minipage}

\newpage

% ==============================================================================
\section{Cấu trúc cơ khí và Hệ thống truyền động Drone X500}
% ==============================================================================

Mẫu Quadrotor được kế thừa từ chuẩn \textbf{PX4 X500} gồm 2 tệp: \path{simulation/models/x500_base/model.sdf} (khung thân, quán tính, 4 cánh quạt, 4 cảm biến) và \path{simulation/models/x500/model.sdf} (4 plugin động cơ + plugin odometry).

% --- 2.1 Base Link ---
\subsection{Khung thân chính (\texttt{base\_link}) và thông số quán tính}
\noindent
\begin{minipage}[t]{0.45\textwidth}
\vspace{0pt}
\textbf{Giải thích cấu hình:}
\begin{itemize}
    \item \textbf{Khối lượng thân:} $m_{\text{body}} = 2.0\text{ kg}$. Cộng thêm $4$ cụm rotor ($4 \times 0.016077\text{ kg}$), tổng khối lượng toàn máy bay là $m = 2.0643\text{ kg}$ (khai báo \texttt{MASS\_KG = 2.064f} trong firmware ESP32).
    \item \textbf{Ten-xơ quán tính:} Đối xứng trục $X$ và $Y$ với $I_{xx} = I_{yy} = 0.02167\text{ kg}\cdot\text{m}^2$, trục quay Yaw $I_{zz} = 0.040\text{ kg}\cdot\text{m}^2$.
    \item \textbf{Tương tác gió:} Bật \texttt{<enable\_wind>true</enable\_wind>} để plugin gió tác động lực đẩy ngang lên thân drone.
\end{itemize}
\end{minipage}
\hfill
\begin{minipage}[t]{0.52\textwidth}
\vspace{0pt}
\begin{lstlisting}[language=SDF, firstnumber=7]
<link name="base_link">
  <inertial>
    <mass>2.0</mass>
    <inertia>
      <ixx>0.02166666666666667</ixx>
      <ixy>0</ixy>
      <ixz>0</ixz>
      <iyy>0.02166666666666667</iyy>
      <iyz>0</iyz>
      <izz>0.04000000000000001</izz>
    </inertia>
  </inertial>
  <gravity>true</gravity>
  <enable_wind>true</enable_wind>
  <velocity_decay />
\end{lstlisting}
\end{minipage}

\vspace{1em}

% --- 2.2 Rotors & Motors ---
\subsection{Cụm 4 Động cơ không chổi than và Cánh quạt 1345 (\texttt{MulticopterMotorModel})}
\noindent
\begin{minipage}[t]{0.45\textwidth}
\vspace{0pt}
\textbf{Giải thích cấu hình:}
\begin{itemize}
    \item \textbf{Bố trí hình học chữ X:} 4 động cơ đặt đối xứng ở $(\pm 0.174,\, \pm 0.174,\, 0.06)\text{ m}$ (cánh tay đòn $L = 0.174\text{ m}$). Rotor 0, 1 quay ngược chiều kim đồng hồ (\texttt{ccw}); Rotor 2, 3 quay thuận chiều (\texttt{cw}).
    \item \textbf{Độ trễ động cơ (Motor Lag):} Hằng số thời gian tăng tốc $\tau_{\text{up}} = 0.0125\text{ s}$, giảm tốc $\tau_{\text{down}} = 0.025\text{ s}$.
    \item \textbf{Hệ số khí động:} Lực nâng $F = k_F \omega^2$ với $k_F = 8.54858 \times 10^{-6}\text{ N}/(\text{rad/s})^2$; mô-men phản lực $M = k_M F$ với $k_M = 0.016\text{ m}$.
    \item \textbf{Tốc độ cực đại:} $\omega_{\max} = 1000\text{ rad/s}$ $\Rightarrow$ lực nâng tối đa 4 động cơ là $4 \times k_F \times 1000^2 \approx 34.194\text{ N}$.
    \item \textbf{Giao tiếp điều khiển:} Nhận lệnh tốc độ \texttt{w0..w3} qua topic \texttt{command/motor\_speed}.
\end{itemize}
\end{minipage}
\hfill
\begin{minipage}[t]{0.52\textwidth}
\vspace{0pt}
\begin{lstlisting}[language=SDF, firstnumber=341]
<!-- Trong x500_base/model.sdf -->
<link name="rotor_0">
  <pose>0.174 -0.174 0.06 0 0 0</pose>
  <inertial>
    <mass>0.016076923076923075</mass>
    <inertia>
      <ixx>3.84649e-07</ixx>
      <iyy>2.61159e-05</iyy>
      <izz>2.64986e-05</izz>
    </inertia>
  </inertial>
</link>
\end{lstlisting}
\vspace{0.4em}
\begin{lstlisting}[language=SDF, firstnumber=7]
<!-- Trong x500/model.sdf (lap lai cho motor 0..3) -->
<plugin filename="gz-sim-multicopter-motor-model-system"
  name="gz::sim::systems::MulticopterMotorModel">
  <jointName>rotor_0_joint</jointName>
  <linkName>rotor_0</linkName>
  <turningDirection>ccw</turningDirection>
  <timeConstantUp>0.0125</timeConstantUp>
  <timeConstantDown>0.025</timeConstantDown>
  <maxRotVelocity>1000.0</maxRotVelocity>
  <motorConstant>8.54858e-06</motorConstant>
  <momentConstant>0.016</momentConstant>
  <commandSubTopic>command/motor_speed</commandSubTopic>
  <motorNumber>0</motorNumber>
  <rotorDragCoefficient>8.06428e-05</rotorDragCoefficient>
  <rollingMomentCoefficient>1e-06</rollingMomentCoefficient>
  <rotorVelocitySlowdownSim>40</rotorVelocitySlowdownSim>
  <motorType>velocity</motorType>
</plugin>
\end{lstlisting}
\end{minipage}

\newpage

% ==============================================================================
\section{Các linh kiện cảm biến trên Drone và thông số nhiễu (\texttt{x500\_base/model.sdf})}
% ==============================================================================

Trên \texttt{base\_link} của Drone X500 gắn \textbf{4 linh kiện cảm biến riêng biệt}. Tất cả đều đặt \texttt{<always\_on>1</always\_on>} để tự động publish liên tục theo đúng \texttt{<update\_rate>} ngay khi Gazebo chạy.

% --- 3.1 IMU ---
\subsection{Cảm biến quán tính IMU (\texttt{imu\_sensor} — 250 Hz)}
\noindent
\begin{minipage}[t]{0.45\textwidth}
\vspace{0pt}
\textbf{Giải thích linh kiện \& thông số:}
\begin{itemize}
    \item \textbf{Chức năng:} Tích hợp 2 bộ đo trong 1 linh kiện: \textbf{Con quay hồi chuyển (Gyroscope)} đo vận tốc góc $(g_x, g_y, g_z)$ và \textbf{Gia tốc kế (Accelerometer)} đo gia tốc $(a_x, a_y, a_z)$.
    \item \textbf{Tần số lấy mẫu:} \texttt{250 Hz} (mỗi $4\text{ ms}$, đồng bộ với step vật lý).
    \item \textbf{Nhiễu Gaussian Gyro:} \texttt{mean = 0.0}, \texttt{stddev = 0.0035 rad/s} ($\approx 0.2^\circ/\text{s}$) trên cả 3 trục.
    \item \textbf{Nhiễu Gaussian Accel:} \texttt{mean = 0.0}, \texttt{stddev = 0.035 m/s$^2$} ($\approx 3.5\text{ mg-rms}$) trên cả 3 trục.
    \item \textbf{Topic GZ-Transport:} \path{/world/default/model/x500/link/base_link/sensor/imu_sensor/imu}.
\end{itemize}
\end{minipage}
\hfill
\begin{minipage}[t]{0.52\textwidth}
\vspace{0pt}
\begin{lstlisting}[language=SDF, firstnumber=255]
<sensor name="imu_sensor" type="imu">
  <gz_frame_id>base_link</gz_frame_id>
  <always_on>1</always_on>
  <update_rate>250</update_rate>
  <imu>
    <angular_velocity>
      <x><noise type="gaussian"><mean>0.0</mean><stddev>0.0035</stddev></noise></x>
      <y><noise type="gaussian"><mean>0.0</mean><stddev>0.0035</stddev></noise></y>
      <z><noise type="gaussian"><mean>0.0</mean><stddev>0.0035</stddev></noise></z>
    </angular_velocity>
    <linear_acceleration>
      <x><noise type="gaussian"><mean>0.0</mean><stddev>0.035</stddev></noise></x>
      <y><noise type="gaussian"><mean>0.0</mean><stddev>0.035</stddev></noise></y>
      <z><noise type="gaussian"><mean>0.0</mean><stddev>0.035</stddev></noise></z>
    </linear_acceleration>
  </imu>
</sensor>
\end{lstlisting}
\end{minipage}

\vspace{1em}

% --- 3.2 GPS / NavSat ---
\subsection{Cảm biến định vị vệ tinh GPS / NavSat (\texttt{navsat\_sensor} — 30 Hz)}
\noindent
\begin{minipage}[t]{0.45\textwidth}
\vspace{0pt}
\textbf{Giải thích linh kiện \& thông số:}
\begin{itemize}
    \item \textbf{Chức năng:} Đo tọa độ địa lý (\texttt{lat, lon, alt}) và vận tốc 3 trục (\texttt{vel\_E, vel\_N, vel\_U}).
    \item \textbf{Tần số lấy mẫu:} \texttt{30 Hz} (chu kỳ $\approx 33.3\text{ ms}$).
    \item \textbf{Nhiễu vị trí ngang:} Trong Gazebo Harmonic, đơn vị \texttt{horizontal} tính theo \textbf{độ}: \texttt{stddev = 4.5e-6 deg} ($\approx 4.5\times 10^{-6} \times 111{,}320 \approx 0.50\text{ m}$).
    \item \textbf{Nhiễu vị trí đứng:} \texttt{stddev = 0.50 m}.
    \item \textbf{Nhiễu vận tốc:} Ngang \texttt{0.15 m/s}; đứng \texttt{0.20 m/s}.
    \item \textbf{Topic GZ-Transport:} \path{/world/default/model/x500/link/base_link/sensor/navsat_sensor/navsat}.
\end{itemize}
\end{minipage}
\hfill
\begin{minipage}[t]{0.52\textwidth}
\vspace{0pt}
\begin{lstlisting}[language=SDF, firstnumber=304]
<sensor name="navsat_sensor" type="navsat">
  <gz_frame_id>base_link</gz_frame_id>
  <always_on>1</always_on>
  <update_rate>30</update_rate>
  <navsat>
    <position_sensing>
      <horizontal>
        <noise type="gaussian"><mean>0.0</mean><stddev>4.5e-6</stddev></noise>
      </horizontal>
      <vertical>
        <noise type="gaussian"><mean>0.0</mean><stddev>0.50</stddev></noise>
      </vertical>
    </position_sensing>
    <velocity_sensing>
      <horizontal>
        <noise type="gaussian"><mean>0.0</mean><stddev>0.15</stddev></noise>
      </horizontal>
      <vertical>
        <noise type="gaussian"><mean>0.0</mean><stddev>0.20</stddev></noise>
      </vertical>
    </velocity_sensing>
  </navsat>
</sensor>
\end{lstlisting}
\end{minipage}

\vspace{1em}

% --- 3.3 Barometer & Magnetometer ---
\subsection{Cảm biến áp suất Barometer (50 Hz) và La bàn từ trường Magnetometer (100 Hz)}
\noindent
\begin{minipage}[t]{0.45\textwidth}
\vspace{0pt}
\textbf{Giải thích linh kiện \& thông số:}
\begin{itemize}
    \item \textbf{Barometer (\texttt{air\_pressure\_sensor}):} Đo áp suất khí quyển tại \texttt{50 Hz} với nhiễu Gaussian \texttt{stddev = 10.0 Pa} ($\approx 0.85\text{ m}$ độ cao). \texttt{StateEstimator} quy đổi áp suất sang độ cao $Z$ theo công thức chuẩn ICAO sau khi hiệu chuẩn $P_0$ từ 15 mẫu đầu trên mặt đất.
    \item \textbf{Magnetometer (\texttt{magnetometer\_sensor}):} Là linh kiện tách riêng khỏi IMU, đo vector từ trường 3 trục tại \texttt{100 Hz} với nhiễu \texttt{stddev = 2e-6 T} trên mỗi trục để hiệu chỉnh trôi góc hướng \texttt{Yaw}.
\end{itemize}
\end{minipage}
\hfill
\begin{minipage}[t]{0.52\textwidth}
\vspace{0pt}
\begin{lstlisting}[language=SDF, firstnumber=219]
<sensor name="air_pressure_sensor" type="air_pressure">
  <gz_frame_id>base_link</gz_frame_id>
  <always_on>1</always_on>
  <update_rate>50</update_rate>
  <air_pressure>
    <pressure>
      <noise type="gaussian"><mean>0</mean><stddev>10.0</stddev></noise>
    </pressure>
  </air_pressure>
</sensor>

<sensor name="magnetometer_sensor" type="magnetometer">
  <gz_frame_id>base_link</gz_frame_id>
  <always_on>1</always_on>
  <update_rate>100</update_rate>
  <magnetometer>
    <x><noise type="gaussian"><stddev>2e-6</stddev></noise></x>
    <y><noise type="gaussian"><stddev>2e-6</stddev></noise></y>
    <z><noise type="gaussian"><stddev>2e-6</stddev></noise></z>
  </magnetometer>
</sensor>
\end{lstlisting}
\end{minipage}

\vspace{1em}

% --- 3.4 Firmware Config Mapping ---
\subsection{Bảng tổng hợp và ánh xạ thông số sang Firmware ESP32 (\texttt{flight\_config.h})}
\noindent
\begin{minipage}[t]{0.45\textwidth}
\vspace{0pt}
\textbf{Giải thích ánh xạ phần cứng:}
\begin{itemize}
    \item \textbf{Nhất quán mô hình:} Toàn bộ thông số khối lượng, mô-men quán tính và hệ số động cơ của X500 trong Gazebo đều được khai báo đồng bộ vào namespace \path{FlightConfig::RLCritic} để bộ quan sát nhiễu \textbf{ESO} và \textbf{MotorMixer} trên ESP32 tính toán đúng lực/mô-men vật lý.
    \item \textbf{Quy đổi thang Mixer:} Lực đẩy tối đa $34.194\text{ N}$, mô-men Roll/Pitch tối đa $4 \times 0.174 \times 8.54858 = 5.950\text{ N}\cdot\text{m}$, mô-men Yaw tối đa $4 \times 0.016 \times 8.54858 = 0.5471\text{ N}\cdot\text{m}$.
    \item \textbf{Bù trễ cơ cấu chấp hành:} Khai báo \path{ESO_ACT_LAG_ATT_S = 0.035 s} khớp với trễ động cơ ($\tau_{\text{up/down}}$) và truyền thông UART.
\end{itemize}
\end{minipage}
\hfill
\begin{minipage}[t]{0.52\textwidth}
\vspace{0pt}
\begin{lstlisting}[language=C++, firstnumber=150]
namespace RLCritic {
    // Thong so vat ly Quadrotor X500 (khop model.sdf)
    constexpr float MASS_KG              = 2.064f;
    constexpr float GRAVITY              = 9.80f;
    constexpr float JX                   = 0.02167f;
    constexpr float JY                   = 0.02167f;
    constexpr float JZ                   = 0.04000f;

    // He so quy doi sang MotorMixer [0..1]
    constexpr float MAX_TOTAL_THRUST_N   = 34.194f;
    constexpr float MIXER_ROLL_SCALE_NM  = 5.950f;
    constexpr float MIXER_PITCH_SCALE_NM = 5.950f;
    constexpr float MIXER_YAW_SCALE_NM   = 0.5471f;

    // Tre dap ung goc & motor
    constexpr float ESO_ACT_LAG_POS_S    = 0.10f;
    constexpr float ESO_ACT_LAG_ATT_S    = 0.035f;
}
\end{lstlisting}
\end{minipage}

\vspace{1.2em}
\noindent
\begin{table}[h!]
\centering
\caption{Tổng hợp các linh kiện và tham số cấu hình chính trên Drone X500}
\small
\begin{tabularx}{\textwidth}{@{}llccX@{}}
\toprule
\textbf{Linh kiện} & \textbf{Tệp cấu hình} & \textbf{Tần số} & \textbf{Nhiễu / Giới hạn} & \textbf{Vai trò trong vòng điều khiển} \\
\midrule
Khung thân \texttt{base\_link} & \texttt{x500\_base/model.sdf} & -- & $m = 2.0\text{ kg}$ & Chịu lực nâng, trọng lực, gió động \\
4 Rotor + Motor 5010 & \texttt{x500/model.sdf} & 250 Hz & $\omega_{\max}=1000\text{ rad/s}$ & Tạo lực nâng ($34.19\text{ N}$) \& mô-men quay \\
IMU (\texttt{imu\_sensor}) & \texttt{x500\_base/model.sdf} & 250 Hz & Gyro $0.0035$, Acc $0.035$ & Đo tốc độ xoay $(p,q,r)$ \& gia tốc \\
GPS (\texttt{navsat\_sensor}) & \texttt{x500\_base/model.sdf} & 30 Hz & Ngang $0.50\text{ m}$, Đứng $0.50\text{ m}$ & Đo tọa độ $(lat,lon,alt)$ \& vận tốc \\
Barometer & \texttt{x500\_base/model.sdf} & 50 Hz & $\sigma = 10.0\text{ Pa}$ ($\approx 0.85\text{ m}$) & Đo áp suất quy đổi độ cao $Z$ \\
Magnetometer & \texttt{x500\_base/model.sdf} & 100 Hz & $\sigma = 2\times 10^{-6}\text{ T}$ & Đo từ trường bù trôi góc \texttt{Yaw} \\
\bottomrule
\end{tabularx}
\end{table}

\end{document}
